% =====================================================================
%  第 1 章 局部 Langlands 对应
%  覆盖：PDF p.15–p.22（原书印刷 p.1–p.8）
%  第 1 章章首页 + §1.1 经典理论 全节（1.1.1–1.1.6）
% =====================================================================

\chapter{局部 Langlands 对应}

在本导论章中，我们详细说明“圈群的局部 Langlands 对应”指的是什么。先简要概述局部非阿基米德域上约化群的局部 Langlands 对应；这类域例如 $\F_q((t))$，即有限域 $\F_q$ 上的洛朗幂级数域。我们要理解的，是把 $\F_q$ 换成复数域 $\C$ 之后这一对应的类比物。此时约化群的角色由形式圈群 $G(\C((t)))$ 扮演。我们遵循 [Frenkel and Gaitsgory 2006c]，讨论在此情形下 Langlands 对应的框架应如何改变。这一讨论将自然把我们引向相应仿射 Kac--Moody 代数 $\gghat$ 的表示范畴，其上带有 $G(\C((t)))$ 的作用；这正是本书其余部分要研究的内容。

% [Tutor Note: 背景（p.1）。把 $\F_q$ 换成 $\C$ 不只是换系数域：$\F_q((t))$ 与 $\C((t))$ 都是洛朗幂级数域，但前者的绝对伽罗瓦群是射影极限（含 wild inertia），后者的绝对伽罗瓦群同构于 $\widehat{\Z}$（见第 1.2.1 小节）；同时 $\C$ 上的“伽罗瓦表示”要换成去心圆盘上的局部系统。全书的核心问题正是：在这一替换下，$W'_F$ 的表示该由什么对象取代。]

\begin{keypoints}
    \item \textbf{核心思想：}把局部 Langlands 对应中的 $\F_q$ 换成 $\C$：局部域 $\F_q((t))$ 换成 $\C((t))$，约化群 $G(\F_q((t)))$ 的角色由形式圈群 $G(\C((t)))$ 扮演，对应关系应改写为仿射 Kac--Moody 代数 $\gghat$ 的表示范畴，并带 $G(\C((t)))$ 的作用。
    \item \textbf{数学工具：}$GL_n$ 的光滑表示、Weil--Deligne 群 $W'_F$ 的容许表示、根数据与 Langlands 对偶群 ${}^L G$。
    \item \textbf{关键结论：}$GL_n$ 的局部对应已是定理；一般约化群仍是猜想，且只对应到 L-包；整体情形在函数域上对任意 $n$ 成立。
    \item \textbf{记号提醒：}$\gghat$ 是圈代数 $\fg(\!(t)\!)$ 的仿射 Kac--Moody 代数，${}^L G$ 是 Langlands 对偶群；公式编号沿用原书。
\end{keypoints}

% =====================================================================
%  1.1 经典理论（印刷 p.1–p.8）
% =====================================================================

\section{经典理论}

\begin{keypoints}
    \item \textbf{核心思想：}局部 Langlands 对应把局部域上约化群的光滑表示与该域伽罗瓦群（经 Weil--Deligne 群）的表示联系起来。
    \item \textbf{数学工具：}局部非阿基米德域、同余子群与光滑表示、绝对伽罗瓦群与 Weil 群、根数据与 Langlands 对偶群。
    \item \textbf{关键结论：}$GL_n$ 情形是定理（函数域 [Laumon et al. 1993]；$p$-adic 域 [Harris and Taylor 2001]、[Henniart 2000]）；一般约化群仍为猜想，且对应到 L-包而非单个表示。
    \item \textbf{记号提醒：}$W'_F=W_F\ltimes\C$ 是 Weil--Deligne 群（式 (1.1-1)），${}^L G$ 是 Langlands 对偶群。
\end{keypoints}

局部 Langlands 对应把局部域上约化代数群的光滑表示与该域的伽罗瓦群的表示联系起来。我们首先定义这两类对象，并解释这一对应的主要特征。本节材料只起动机作用，因此我们只提及其中最与我们相关的方面。更详细的处理请读者参阅颇有教益的综述 [Vogan 1993; Kudla 1994] 及其中的参考文献。

\subsection{局部非阿基米德域}

\begin{keypoints}
    \item \textbf{核心思想：}局部非阿基米德域只有两类：$p$-adic 域的有限扩张，以及 $\F_q((t))$ 的有限扩张。
    \item \textbf{关键结论：}$\Q_p$ 与 $\F_q((t))$ 作为拓扑空间都由“从某项起的级数”给出，范数由首个非零项的次数决定。
    \item \textbf{易错点：}$p$-adic 数有“进位”，$\F_q((t))$ 没有；两者特征不同（$0$ 与 $p$），级数形式相似并不表示域同构。
\end{keypoints}

设 $F$ 为局部非阿基米德域。换言之，$F$ 是 $p$-adic 数域 $\Q_p$，或 $\Q_p$ 的有限扩张；或者 $F$ 是系数在 $\F_q$（含 $q$ 个元素的有限域）中的形式洛朗幂级数域 $\F_q((t))$。\transnote{原文对正特征情形只举了 $\F_q((t))$，未写“有限扩张”。实际上正特征局部域都是 $\F_q((t))$ 的有限扩张，总同构于某个 $\F_{q'}((t'))$，与上一类“$\Q_p$ 的有限扩张”正好对称；原书 p.7 对任意闭点写 $F_x\simeq\F_{q_x}((t_x))$，用的正是这一般形状。两类域的特征分别为 $0$ 与 $p$，级数写法相似并不表示二者同构。[依据：原书 p.1 的列举与 p.7 的 $F_x$ 写法。][来源：本书 p.1。]}

我们回忆：对形如 $p^n$ 的任意 $q$（$p$ 为素数），存在唯一（在同构意义下）的特征为 $p$、含 $q$ 个元素的有限域。

% -------------------------- 印刷 p.2（PDF p.16） --------------------------

$\F_q((t))$ 的元素是形如
\[
    \sum_{n\in\Z} a_n t^n,\qquad a_n\in\F_q
\]
的表达式，其中存在整数 $N$，使得 $n<N$ 时 $a_n=0$。换言之，这些级数沿正方向无限、沿负方向有限。回忆 $p$-adic 数也可写成级数
\[
    \sum_{n\in\Z} b_n p^n,\qquad b_n\in\{0,1,\dots,p-1\},
\]
其中存在整数 $N$，使得 $n<N$ 时 $b_n=0$。于是 $\Q_p$ 的元素与 $\F_p((t))$ 的元素看起来相像。两个域都关于下述拓扑完备：该拓扑由范数给出；若级数中 $a_N\neq 0$ 而一切 $n<N$ 的 $a_n=0$，则范数取值 $\alpha^{-N}$，其中 $\alpha$ 是 $0$ 与 $1$ 之间固定的正实数。但两个域的加法与乘法法则不同：$\Q_p$ 情形有“进位”，$\F_p((t))$ 情形没有“进位”。特别地，$\Q_p$ 的特征为 $0$，而 $\F_q((t))$ 的特征为 $p$。更一般地，$\Q_p$ 的有限扩张的元素与某个 $q=p^n$ 对应的 $\F_q((t))$ 的元素看起来相像；但加法与乘法法则以及特征仍然不同。

% [Tutor Note: 易错（p.2）。此处范数取 $\alpha^{-N}$ 而 $\alpha\in(0,1)$，即 $\abs{t}=\alpha^{-1}>1$，这是“赋值即指数”的约定；通常局部域教材把 $\abs{t}$ 归一到 $q^{-1}$。两种约定只差把 $\alpha$ 换成 $\alpha^{-1}$，但下文的 $q^{\nu(\sigma)}$（式 (1.1-1)）与此处的 $\alpha^{-N}$ 必须统一到同一约定，否则指数符号会整体反号。]

\subsection{$GL_n(F)$ 的光滑表示}

\begin{keypoints}
    \item \textbf{核心思想：}光滑表示就是每个向量都被某个开紧子群 $K_N$ 固定的表示。
    \item \textbf{数学工具：}同余子群 $K_N=\{g\in GL_n(\F_q[[t]])\mid g\equiv 1\bmod t^N\}$ 给出 $GL_n(F)$ 的开邻域基。
    \item \textbf{关键结论：}$GL_n(F)$ 的不可约光滑表示，意外地与性质完全不同的对象（$F$ 的伽罗瓦群的 $n$ 维表示）相关。
\end{keypoints}

现在考虑群 $GL_n(F)$，其中 $F$ 是局部非阿基米德域。$GL_n(F)$ 在复向量空间 $V$ 上的表示是满足 $\pi(gh)=\pi(g)\pi(h)$ 与 $\pi(1)=\mathrm{Id}$ 的同态 $\pi:GL_n(F)\to\End V$。在 $GL_n(F)$ 上定义拓扑：规定 $1\in GL_n(F)$ 的开邻域基由同余子群 $K_N$（$N\in\Z_+$）组成。当 $F=\F_q((t))$ 时，群 $K_N$ 定义如下：
\[
    K_N=\{g\in GL_n(\F_q[[t]])\mid g\equiv 1 \bmod t^N\},
\]
而当 $F=\Q_p$ 时，它的定义类似。对每个 $v\in V$ 得到映射 $\pi(\cdot)v:GL_n(F)\to V$，$g\mapsto\pi(g)v$。表示 $(V,\pi)$ 称为\textbf{光滑}的，若对每个 $v$，映射 $\pi(\cdot)v$ 连续，这里 $V$ 赋予离散拓扑。\transnote{这里的“光滑”是表示论意义下的，与 $C^\infty$ 光滑无关。$V$ 赋离散拓扑时，$\pi(\cdot)v$ 连续等价于稳定化子 $\operatorname{Stab}(v)=\{g\in GL_n(F)\mid \pi(g)v=v\}$ 是开子群；又因 $K_N$ 构成 $1$ 的开邻域基，这等价于存在 $N$ 使 $v\in V^{K_N}$（$K_N$-不变向量）。故光滑表示即 $V=\bigcup_N V^{K_N}$，每个向量都被某个紧开子群固定；这是 $p$-adic 表示论的起点，容许性再要求每个 $V^{K_N}$ 有限维。[来源：原书 p.2；并见 [Kudla 1994]，pp. 365--391。]}换言之，$V$ 光滑当且仅当对任意向量 $v\in V$ 存在 $N\in\Z_+$ 使得
\[
    \pi(g)v=v,\qquad \forall g\in K_N.
\]

% [Tutor Note: 易错（p.2，1.1.2）。这里的“光滑”是表示论意义下的光滑（每个向量被某个同余子群 $K_N$ 固定），与 $C^\infty$ 光滑无关；技术上它等价于：$V$ 赋离散拓扑时，稳定化子 $\operatorname{Stab}(v)$ 是开子群。$K_N$ 是 $GL_n(F)$ 的紧开子群，故光滑表示是 $K_N$-有限向量的并，这是后文容许性讨论的起点。]

我们关心的是描述 $GL_n(F)$ 的不可约光滑表示的等价类。出人意料的是，它们与另一类对象相关：$F$ 的伽罗瓦群的 $n$ 维表示。

% -------------------------- 印刷 p.3（PDF p.17） --------------------------

\subsection{伽罗瓦群}

\begin{keypoints}
    \item \textbf{核心思想：}伽罗瓦群只记录可分扩张；取最大可分扩张 $\overline{F}$ 得到绝对伽罗瓦群 $\operatorname{Gal}(\overline{F}/F)$。
    \item \textbf{数学工具：}$\F_q((t))$ 的不可分扩张（如 $\F_q((t^{1/p}))$）不贡献伽罗瓦群；可分扩张 $\F_q((t^{1/n}))$ 仅当 $\zeta_n\in\F_q$（即 $n\mid q-1$）时为伽罗瓦扩张，其伽罗瓦群为 $\Z/n\Z$（正文按原书保留 $(\Z/n\Z)^\times$，见译注）。
    \item \textbf{关键结论：}$\operatorname{Gal}(\overline{\F}_q/\F_q)\simeq\widehat{\Z}=\varprojlim\Z/N\Z$；$\Z\subset\widehat{\Z}$ 的原像给出 Weil 群 $W_F$。
    \item \textbf{记号提醒：}式 (1.1-1) 定义 $W'_F=W_F\ltimes\C$；$W'_F$ 的容许表示是一对 $(\rho,u)$，$\rho$ 半单连续、$u$ 幂零。
\end{keypoints}

设 $F$ 是 $K$ 的子域。则\textbf{伽罗瓦群} $\operatorname{Gal}(K/F)$ 由 $K$ 的所有满足 $\sigma(y)=y$（对一切 $y\in F$）的域自同构 $\sigma$ 组成。

设 $F$ 为域。$F$ 的代数闭包是把系数在 $F$ 中的一切多项式的根添加到 $F$ 上所得到的域。当 $F=\F_q((t))$ 时，$F$ 的某些扩张可能不可分。这种扩张的一个例子是域 $\F_q((t^{1/p}))$。定义该扩张的多项式是 $x^p-t$，但它在 $\F_q((t^{1/p}))$ 中有重根，因为
\[
    x^p-(t^{1/p})^p=(x-t^{1/p})^p.
\]
该扩张的伽罗瓦群 $\operatorname{Gal}(\F_q((t^{1/p})),\F_q((t)))$ 是平凡群，尽管扩张的次数为 $p$。\transnote{原书此处把伽罗瓦群的记号印成了 $\operatorname{Gal}(\F_q((t^{1/p})),\F_q((t)))$，两个域之间是逗号。本页上文与下文都用斜杠记号（如 $\operatorname{Gal}(\overline{F}/F)$），按惯例应写作 $\operatorname{Gal}(\F_q((t^{1/p}))/\F_q((t)))$。译文按原书保留逗号，公式未改。[依据：原书 p.3 同页记号。][来源：本书 p.3。]}

% [Tutor Note: 原书笔误（p.3）。上式伽罗瓦群记号中间是逗号，按惯例应为斜杠，即 $\operatorname{Gal}(\F_q((t^{1/p}))/\F_q((t)))$。译文保留原书逗号，已登记到报告“原书疑似笔误”栏。]

这一扩张应与可分扩张 $\F_q((t^{1/n}))$ 对照，其中 $n$ 不被 $p$ 整除。该扩张由多项式 $x^n-t$ 定义，此时它没有重根：
\[
    x^n-t=\prod_{i=0}^{n-1}(x-\zeta^i t^{1/n}),
\]
其中 $\zeta$ 是 $\F_q$ 的代数闭包中的 $n$ 次本原单位根。相应的伽罗瓦群与群 $(\Z/n\Z)^\times$ 等同，即 $\Z/n\Z$ 的可逆元群。\transnote{原文说相应的伽罗瓦群是 $(\Z/n\Z)^\times$，这与扩张次数为 $n$ 不相容：$x^n-t$ 在 $\F_q((t))$ 上是 Eisenstein 多项式，故 $K=\F_q((t^{1/n}))$ 在 $F=\F_q((t))$ 上的次数为 $n$，而 $(\Z/n\Z)^\times$ 的阶为 $\varphi(n)$。自同构必形如 $t^{1/n}\mapsto\zeta^a t^{1/n}$（$a\in\Z/n\Z$），且它落在 $K$ 的自同构群中当且仅当 $\zeta^a\in\F_q$：故当 $\zeta_n\in\F_q$（即 $n\mid q-1$）时 $K/F$ 是伽罗瓦扩张，群为 $\Z/n\Z$；否则 $K/F$ 不是伽罗瓦扩张，$\operatorname{Aut}(K/F)\cong\mu_n\cap\F_q\cong\Z/\gcd(n,q-1)$。原书的 $(\Z/n\Z)^\times$ 疑为把分圆扩张的伽罗瓦群误置于此，公式按原书保留。[译者按（无外部依据）：以上由 $x^n-t$ 的根及 $K$ 的剩余域仍为 $\F_q$ 直接验证。]}

我们希望避开不可分扩张，因为它们对伽罗瓦群没有贡献。（若 $F$ 的特征为 $0$，例如 $F=\Q_p$，则不存在不可分扩张。）设 $\overline{F}$ 是给定 $F$ 的代数闭包中的最大可分扩张。它在同构意义下由 $F$ 唯一确定。

称 $\operatorname{Gal}(\overline{F}/F)$ 为 $F$ 的\textbf{绝对伽罗瓦群}。它的元素是 $\overline{F}$ 的满足 $\sigma(y)=y$（对一切 $y\in F$）的域自同构 $\sigma$。\transnote{只取最大可分扩张，是因为不可分扩张不贡献自同构。对 $F=\F_q((t))$，$\operatorname{Gal}(\overline{F}/F)$ 含有 wild inertia（$p$-部分），是相当复杂的射影极限；而对 $F=\C((t))$，代数闭包就是 Puiseux 级数域 $\bigcup_{n\ge1}\C((t^{1/n}))$，其伽罗瓦群同构于 $\widehat{\Z}$（见第 1.2.1 小节）。正是这一差别使几何情形中伽罗瓦一侧大为简化。[依据：原书 p.3；$\C((t))$ 的情形见第 1.2.1 小节。][来源：本书 p.3。]}

% [Tutor Note: 背景（p.3）。只取最大可分扩张，是因为不可分扩张不贡献伽罗瓦群；对 $F=\F_q((t))$，$\operatorname{Gal}(\overline{F}/F)$ 含 wild inertia（$p$-部分），而对 $F=\C((t))$ 它同构于 $\widehat{\Z}$（第 1.2.1 小节）。这一差别正是几何情形中伽罗瓦一侧大为简化的根源。]

为获得一些关于伽罗瓦群的经验，我们来看伽罗瓦群 $\operatorname{Gal}(\overline{\F}_q/\F_q)$。这里 $\overline{\F}_q$ 是 $\F_q$ 的代数闭包，可定义为域 $\F_{q^N}$（$N\in\Z_+$）关于自然嵌入 $\F_{q^N}\hookrightarrow\F_{q^M}$（$N$ 整除 $M$）的归纳极限。因此 $\operatorname{Gal}(\overline{\F}_q/\F_q)$ 同构于伽罗瓦群 $\operatorname{Gal}(\F_{q^N}/\F_q)$ 关于自然满射
\[
    \operatorname{Gal}(\F_{q^M}/\F_q)\twoheadrightarrow\operatorname{Gal}(\F_{q^N}/\F_q),
    \qquad \forall\, N\mid M
\]
的逆向极限。群 $\operatorname{Gal}(\F_{q^N}/\F_q)$ 容易描述：它由 \textbf{Frobenius 自同构} $x\mapsto x^q$ 生成（注意它保持 $\F_q$ 不变），该自同构的阶为 $N$，故 $\operatorname{Gal}(\F_{q^N}/\F_q)\simeq\Z/N\Z$。于是得到
\[
    \operatorname{Gal}(\overline{\F}_q/\F_q)\simeq\widehat{\Z}
    \overset{\mathrm{def}}{=}\varprojlim\Z/N\Z,
\]

% -------------------------- 印刷 p.4（PDF p.18） --------------------------

其中对 $N\mid M$ 有满射 $\Z/M\Z\twoheadrightarrow\Z/N\Z$。群 $\widehat{\Z}$ 包含 $\Z$ 作为子群。

设 $F=\F_q((t))$。注意有自然映射 $\operatorname{Gal}(\overline{F}/F)\to\operatorname{Gal}(\overline{\F}_q/\F_q)$，它由把 $\overline{F}$ 的自同构限制到子域 $\overline{\F}_q\subset\overline{F}$ 上得到。当 $F$ 的特征为 $0$ 时也存在类似的映射。设 $W_F$ 是子群 $\Z\subset\operatorname{Gal}(\overline{\F}_q/\F_q)$ 在该映射下的原像。这就是 $F$ 的 \textbf{Weil 群}。设 $\nu$ 是相应的同态 $W_F\to\Z$。设 $W'_F=W_F\ltimes\C$ 是 $W_F$ 与一维复加法群 $\C$ 的半直积，其中 $W_F$ 通过公式
\begin{equation}
    \sigma x\sigma^{-1}=q^{\nu(\sigma)}x,\qquad \sigma\in W_F,\; x\in\C
    \tag{1.1-1}
\end{equation}
作用在 $\C$ 上。这就是 $F$ 的 \textbf{Weil--Deligne 群}。\transnote{式 (1.1-1) 中 $W_F$ 通过 $\nu$ 作用在加法群 $\C$ 上，特征由 $\sigma\mapsto q^{\nu(\sigma)}$ 给出，所以半直积 $W'_F=W_F\ltimes\C$ 的乘法完全由 $\nu$ 与 $q$ 决定。一个表示 $\rho'$ 由 $\rho=\rho'|_{W_F}$ 与 $u=\mathrm{d}\rho'(1)\in\fgl_n(\C)$ 决定（$\C$ 的像必幂零），而 $\rho'$ 是同态当且仅当 $\rho(\sigma)u\rho(\sigma)^{-1}=q^{\nu(\sigma)}u$ 对一切 $\sigma\in W_F$ 成立——这正是式 (1.1-1) 的像。反过来，任一满足该关系的 $(\rho,u)$ 通过 $\rho'(\sigma x)=\rho(\sigma)\exp(xu)$ 给出 $W'_F$ 的表示，故两种描述等价。[来源：原书 p.4；并见 [Deligne 1973]，pp. 501--597。]}

$W'_F$ 的 $n$ 维复表示按定义是同态 $\rho':W'_F\to GL_n(\C)$，它可描述为一对 $(\rho,u)$：其中 $\rho$ 是 $W_F$ 的 $n$ 维表示，$u\in\fgl_n(\C)$，且对所有 $\sigma\in W_F$ 有 $\rho(\sigma)u\rho(\sigma)^{-1}=q^{\nu(\sigma)}u$。群 $W_F$ 是拓扑群，赋予 Krull 拓扑（其中单位元的开邻域是有限指数的正规子群）。表示 $(\rho,u)$ 称为\textbf{容许}的，若 $\rho$ 连续（等价地，通过 $W_F$ 的某个有限商因子分解）且半单，并且 $u$ 是 $\fgl_n(\C)$ 的幂零元。

% [Tutor Note: 易错（p.4，式 (1.1-1)）。$W_F$ 是 $\Z\subset\operatorname{Gal}(\overline{\F}_q/\F_q)$ 在该自然映射下的原像，所以它不是闭子群；在 Krull 拓扑下 $\nu:W_F\to\Z$ 并不连续（$\Z$ 取离散拓扑）。半直积 $W'_F=W_F\ltimes\C$ 中 $\C$ 是加法群，作用是乘以 $q^{\nu(\sigma)}$；这正解释了 $\rho(\sigma)u\rho(\sigma)^{-1}=q^{\nu(\sigma)}u$ 中 $u$ 的“权”。$\nu$ 在部分文献中记为 $v$ 或 $\operatorname{ord}$。]

群 $W'_F$ 由 P. Deligne [1973] 引入。其想法是：把幂零元 $u$ 添加到 $W_F$ 上，所得到的群的复容许表示与 $W_F$ 的连续 $\ell$-adic 表示相同（其中 $\ell\neq p$ 为素数）。

% [Tutor Note: 含糊点拨（p.4）。Deligne 说两类表示“相同”，准确含义是：由 Grothendieck 单值化定理，$W_F$ 的连续 $\ell$-adic 表示可分解为半单部分 $\rho$ 与幂零单值性 $u$，于是两类对象在等价类层面有自然双射。这并不意味着两个张量范畴等价（$\ell$-adic 一侧还带 Frobenius 的权、$\ell$ 的选择等数据），且必须 $\ell\neq p$。]

\subsection{$GL_n$ 的局部 Langlands 对应}

\begin{keypoints}
    \item \textbf{核心思想：}局部 Langlands 对应是 $GL_n(F)$ 的不可约光滑表示与 $W'_F$ 的 $n$ 维容许表示之间的双射。
    \item \textbf{数学工具：}$W'_F$ 的容许表示即半单连续 $\rho$ 加上幂零 $u$；这是“伽罗瓦一侧”的全部数据。
    \item \textbf{关键结论：}该对应已证：$F=\F_q((t))$ 见 [Laumon et al. 1993]；$F=\Q_p$ 及其有限扩张见 [Harris and Taylor 2001] 与 [Henniart 2000]。
    \item \textbf{易错点：}对应还须满足一组超定的约束（$\varepsilon$-因子、$L$-因子、$\gamma$-因子的匹配等），仅“存在双射”远不足以刻画它。
\end{keypoints}

现在可以陈述局部非阿基米德域 $F$ 上群 $GL_n$ 的局部 Langlands 对应了。它是两类不同数据之间的双射。一类是 $GL_n(F)$ 的不可约光滑表示的等价类集合。另一类是 $W'_F$ 的 $n$ 维容许表示的等价类集合。我们示意如下：
\begin{center}
\fbox{\parbox[c]{0.36\linewidth}{\centering $n$ 维容许表示\\ 的 $W'_F$}}
\quad$\Longleftrightarrow$\quad
\fbox{\parbox[c]{0.36\linewidth}{\centering 不可约光滑表示\\ 的 $GL_n(F)$}}
\end{center}

该对应应当满足一组超定的约束条件，这里我们不再追述（例如见 [Kudla 1994]）。

$GL_n$ 的局部 Langlands 对应是一个定理。当 $F=\F_q((t))$ 时，它已在 [Laumon et al. 1993] 中证明；当 $F=\Q_p$ 或其有限扩张时，已在 [Harris and Taylor 2001] 以及 [Henniart 2000] 中证明。详细内容请读者参阅这些论文以及综述 [Carayol 2000]。

% -------------------------- 印刷 p.5（PDF p.19） --------------------------

尽管过去二十年为理解它投入了巨大努力，局部 Langlands 对应仍然是个谜。我们已知上述双射存在，却还不能完全满意地解释它为何存在。我们不知道使这样的对应成为可能的深层原因。理解它的一个途径是看它有多一般。下一节我们将讨论这一对应的一种可能推广：把群 $GL_n$ 换成 $F$ 上任意的约化代数群。

% [Tutor Note: 背景（p.5）。这一段是全书动机的枢纽：对应“存在但不知为何存在”，正是要寻找一个能解释它的框架。Frenkel 的答案（第 10 章）是把离散群 $G(F)$ 换成圈群 $G(\!(t)\!)$、把向量空间上的表示换成范畴上的表示，使对应成为“中心同构于 opers 上函数”这一代数事实的几何翻译。]

\subsection{到其他约化群的推广}

\begin{keypoints}
    \item \textbf{核心思想：}一般约化群的表示应由 $W'_F$ 到 Langlands 对偶群 ${}^L G$ 的容许同态参数化，而不只是到 $GL_n(\C)$ 的同态。
    \item \textbf{数学工具：}根数据 $(X^*(T),X_*(T),\Delta,\Delta^\vee)$；交换权格与余权格、单根与余根，得到 ${}^L G^\circ$，再与伽罗瓦群作半直积得到 ${}^L G$。
    \item \textbf{关键结论：}一般情形比 $GL_n$ 微妙：须同时考虑所有内形式，且一个同态对应一个 L-包而非单个表示。
\end{keypoints}

让我们把群 $GL_n$ 换成局部非阿基米德域 $F$ 上任意的连通约化群 $G$。群 $G(F)$ 也是拓扑群，并且有 $G(F)$ 在复向量空间上的光滑表示的概念。一个自然的问题是：能否把 $G(F)$ 的不可约光滑表示与 Weil--Deligne 群 $W'_F$ 的表示联系起来？这个问题由一般的局部 Langlands 猜想处理。在此给出这些猜想的精确表述会离题太远。因此我们只指出其中涉及的某些对象，并请读者参阅文章 [Vogan 1993; Kudla 1994]，那里对这些猜想有很详细的描述。

回忆当 $G=GL_n$ 时，不可约光滑表示由容许同态 $W'_F\to GL_n(\C)$ 参数化。对一般的约化群 $G$，猜测表示由 $W'_F$ 到所谓 Langlands 对偶群 ${}^L G$ 的容许同态参数化，而 ${}^L G$ 定义在 $\C$ 上。

为解释 Langlands 对偶群的概念，先考虑域 $F$ 的闭包 $\overline{F}$ 上的群 $G$。该群的所有极大环面 $T$ 彼此共轭，且必然分裂，即同构 $T(\overline{F})\simeq(\overline{F}^\times)^n$。例如当 $G=GL_n$ 时，所有极大环面都共轭于对角矩阵构成的子群。我们给 $T(\overline{F})$ 配两个格：同态 $T(\overline{F})\to\overline{F}^\times$ 构成的权格 $X^*(T)$，以及同态 $\overline{F}^\times\to T(\overline{F})$ 构成的余权格 $X_*(T)$。它们分别包含根系 $\Delta\subset X^*(T)$ 与余根系 $\Delta^\vee\subset X_*(T)$。四元组 $(X^*(T),X_*(T),\Delta,\Delta^\vee)$ 称为 $G$ 在 $\overline{F}$ 上的根数据（root datum）。根数据把 $G$ 确定到差一个定义在 $\overline{F}$ 上的同构。取包含 $T(\overline{F})$ 的 Borel 子群 $B(\overline{F})$，等价于在 $\Delta$ 中取一组基，即单根系 $\Delta_s$，以及 $\Delta^\vee$ 中相应的基 $\Delta^\vee_s$。

% [Tutor Note: 易错（p.5，1.1.5）。权格与余权格的方向极易写反：$X^*(T)=\Hom(T(\overline{F}),\overline{F}^\times)$，$X_*(T)=\Hom(\overline{F}^\times,T(\overline{F}))$，$\Delta\subset X^*(T)$，$\Delta^\vee\subset X_*(T)$。做对偶时要同时交换 $X^*$ 与 $X_*$、$\Delta_s$ 与 $\Delta^\vee_s$（见 p.6），得到的才是 ${}^L G^\circ$ 的带基根数据；只换格不换根会得到错误的群。]

现在给定 $\gamma\in\operatorname{Gal}(\overline{F}/F)$，存在 $g\in G(\overline{F})$ 使得 $g(\gamma(T(\overline{F})))g^{-1}=T(\overline{F})$ 且 $g(\gamma(B(\overline{F})))g^{-1}=B(\overline{F})$。于是 $g$ 给出带基根数据（based root data）$(X^*(T),X_*(T),\Delta_s,\Delta^\vee_s)$ 的一个自同构。这样便得到 $\operatorname{Gal}(\overline{F}/F)$ 在带基根数据上的作用。

% -------------------------- 印刷 p.6（PDF p.20） --------------------------

现在我们交换权格与余权格，并交换单根与余根。于是得到带基根数据
\[
    (X_*(T),X^*(T),\Delta^\vee_s,\Delta_s),
\]
它是某个 $\C$ 上约化代数群的带基根数据，该群记为 ${}^L G^\circ$。例如，群 $GL_n$ 自对偶，$SO_{2n+1}$ 的对偶是 $Sp_{2n}$，$Sp_{2n}$ 的对偶是 $SO_{2n+1}$，而 $SO_{2n}$ 自对偶。\transnote{对偶的不是群本身，而是带基根数据：$(X^*(T),X_*(T),\Delta,\Delta^\vee)$ 记录权格、余权格与根、余根，选定 Borel 子群相当于选定单根系 $\Delta_s$，得带基根数据。同时交换两个格与 $\Delta_s,\Delta^\vee_s$，所得仍是某 $\C$ 上约化群的带基根数据（Chevalley 定理），记该群为 ${}^L G^\circ$；于是 $GL_n$ 自对偶、$SO_{2n+1}$ 与 $Sp_{2n}$ 互换（$B_n\leftrightarrow C_n$）、$SO_{2n}$ 自对偶（$D_n$）。$\operatorname{Gal}(\overline{F}/F)$ 经其在带基根数据上的作用作用于 ${}^L G^\circ$，半直积 ${}^L G=\operatorname{Gal}(\overline{F}/F)\ltimes{}^L G^\circ$ 才是完整的 Langlands 对偶群；它一般不是 $\C$ 上的代数群，Langlands 本人称其为 L-群（L-group）。[来源：原书 p.5--6；L-群之名见 nLab “Langlands dual group” 的 Remark on terminology（访问 2026-11-23）。]}

$\operatorname{Gal}(\overline{F}/F)$ 在带基根数据上的作用诱导出它在 ${}^L G^\circ$ 上的作用。半直积 ${}^L G=\operatorname{Gal}(\overline{F}/F)\ltimes{}^L G^\circ$ 称为 $G$ 的 \textbf{Langlands 对偶群}。

% [Tutor Note: 易错（p.6）。对偶按根系类型成对出现：$B_n$ 与 $C_n$ 互换、$D_n$ 自对偶，故 $SO_{2n+1}$ 与 $Sp_{2n}$ 互换、$SO_{2n}$ 自对偶。但 ${}^L G^\circ$ 由完整的带基根数据决定（权格与余权格的交换），只凭根系类型判断在一般约化群上会出错；$GL_n$ 自对偶来自其根数据的自对偶。]

群 $G(F)$ 的局部 Langlands 对应把 $G(F)$ 的不可约光滑表示的等价类与容许同态 $W'_F\to{}^L G$ 的等价类联系起来。不过，一般而言这一对应比 $GL_n$ 情形微妙得多。特别地，我们需要同时考虑 $G$ 的所有内形式的表示；并且同态 $W'_F\to{}^L G$ 一般不对应 $G(F)$ 的单个不可约表示，而对应一个有限表示集，称为 \textbf{L-包}（$L$-packet）。\transnote{为何一般群只能得到 L-包：$G$ 的所有内形式共享同一个 ${}^L G$，局部对应必须同时覆盖它们；同一参数 $W'_F\to{}^L G$ 下的不可约表示构成有限集，包内元素由内窥（endoscopy）与稳定分布理论相联系，需附加数据才能分离。$GL_n$ 的 L-包恰为单点，所以 $GL_n$ 的情形看起来像双射；一般群的对应应读作“参数 $\mapsto$ L-包”。[来源：原书 p.6 与参考文献 [Vogan 1993]，pp. 305--379；例子见本书第 10.4.1 节。]}为区分它们，需要额外的数据（详见 [Vogan 1993]；一些例子见下面第 10.4.1 节）。但作为初步近似可以说，局部 Langlands 对应的实质是：
\begin{center}
\itshape 群 $G(F)$ 的不可约光滑表示由容许同态 $W'_F\to{}^L G$ 参数化。
\end{center}

% [Tutor Note: 含糊点拨（p.6）。为何一般约化群只能得到 L-包：同一个 $W'_F\to{}^L G$ 无法区分 L-包内的表示，这些表示由内窥（endoscopy）与稳定分布理论相联系，需要附加数据来分离（见 Vogan 1993）。一个具体来源是“同时考虑所有内形式”：例如 $GL_n$ 的内形式是 $GL_m(D)$（$D$ 为 $F$ 上中心可除代数，$md=n$），$GL_n$ 情形之所以干净，是因为其 L-包恰为单点。]

\subsection{关于整体 Langlands 对应}

\begin{keypoints}
    \item \textbf{核心思想：}整体对应在函数域上把 $W_F$ 的 $n$ 维表示与 $GL_n(\mathbb{A})$ 的不可约自守表示配对；它应在每个点局部化后与局部对应相容。
    \item \textbf{数学工具：}曲线 $X$ 的闭点给出局部域 $F_x$；限制乘积给出阿代尔环 $\mathbb{A}=\mathbb{A}_F$，限制张量积给出整体表示。
    \item \textbf{关键结论：}全局对应对 $GL_2$（Drinfeld）与任意 $n$（Lafforgue）已证；数域情形只是单射，不是双射。
    \item \textbf{易错点：}$F$ 经对角嵌入成为 $\mathbb{A}$ 的离散子环；商 $GL_n(F)\backslash GL_n(\mathbb{A})$ 才承载自守函数。
\end{keypoints}

本节以对整体 Langlands 对应及其与局部对应之联系的简要讨论作结。我们将在第 10.5 节回到这一主题。

设 $X$ 是 $\F_q$ 上的光滑射影曲线。记 $F$ 为 $X$ 上的有理函数域 $\F_q(X)$。对 $X$ 的任一闭点 $x$，记 $F_x$ 为 $F$ 在 $x$ 处的完备化，记 $\cO_x$ 为其整数环。若在 $x$ 处取局部坐标 $t_x$（即 $X$ 上在 $x$ 处一阶消失的有理函数），则得到同构 $F_x\simeq\F_{q_x}((t_x))$ 与 $\cO_x\simeq\F_{q_x}[[t_x]]$，其中 $\F_{q_x}$ 是 $x$ 的剩余域；一般地，它是 $\F_q$ 的有限扩张，含 $q_x=q^{\deg(x)}$ 个元素。

% -------------------------- 印刷 p.7（PDF p.21） --------------------------

于是，现在 $X$ 的每个点上都附有一个局部域。$F$ 的阿代尔环 $\mathbb{A}=\mathbb{A}_F$ 按定义是域 $F_x$ 的限制乘积，其中 $x$ 遍历 $X$ 的所有闭点构成的集合 $\abs{X}$。“限制”一词的意思是：我们只考虑由 $F_x$ 的元素组成的族 $(f_x)_{x\in\abs{X}}$，其中除有限多个 $x$ 外都有 $f_x\in\cO_x$。环 $\mathbb{A}$ 包含域 $F$；$F$ 通过取有理函数在 $X$ 各点的展开，对角地嵌入 $\mathbb{A}$。\transnote{若改用直积 $\prod_x F_x$，$F$ 的对角嵌入将不再是离散的，商空间也失去好的测度性质；限制乘积只允许除有限多个 $x$ 外 $f_x\in\cO_x$，正是为此。表示一侧的对应物是限制张量积 $\bigotimes'_{x}\pi_x$：除有限多个 $x$ 外 $\pi_x$ 无分歧，取其一维 $K_x$-不变向量作单位向量，张量积才有定义。函数域情形还能进一步调用曲线 $X$ 与其模空间的几何，这正是整体对应在函数域率先被证明的原因之一（见 p.8）。[来源：原书 p.7；并见第 10.5 节。]}

% [Tutor Note: 易错（p.7）。“限制乘积”的“限制”指：除有限多个 $x$ 外 $f_x\in\cO_x$；对应到表示一侧，整体自守表示的局部分量除有限多个 $x$ 外都是无分歧的，这正是限制张量积 $\bigotimes'_{x}\pi_x$ 有意义的原因。另外，$F$ 经对角嵌入成为 $\mathbb{A}$ 的离散子环（函数域情形下还是余紧的），否则商 $GL_n(F)\backslash GL_n(\mathbb{A})$ 上无法谈不变测度与自守形式。]

局部 Langlands 对应中考虑的是局部域上群 $GL_n$ 的不可约光滑表示，而整体 Langlands 对应中考虑的是群 $GL_n(\mathbb{A})$ 的不可约\textbf{自守表示}。“自守”一词大致是说，该表示可以实现在商 $GL_n(F)\backslash GL_n(\mathbb{A})$ 上某个合理的函数空间里（$GL_n(\mathbb{A})$ 从右作用）。

对应的另一侧，我们考虑伽罗瓦群 $\operatorname{Gal}(\overline{F}/F)$ 的 $n$ 维表示；更确切地说，是 Weil 群 $W_F$ 的表示。$W_F$ 是 $\operatorname{Gal}(\overline{F}/F)$ 的子群，其定义与局部情形相同。

粗略地说，整体 Langlands 对应是 $W_F$ 的 $n$ 维表示的等价类集合与 $GL_n(\mathbb{A})$ 的不可约自守表示的等价类集合之间的双射：
\begin{center}
\fbox{\parbox[c]{0.36\linewidth}{\centering $n$ 维表示\\ 的 $W_F$}}
\quad$\Longleftrightarrow$\quad
\fbox{\parbox[c]{0.40\linewidth}{\centering 不可约自守表示\\ 的 $GL_n(\mathbb{A})$}}
\end{center}

精确陈述要微妙些。例如，我们应考虑 Weil 群的所谓 $\ell$-adic 表示（而在局部情形考虑的是 Weil--Deligne 群的复容许表示；原因是局部情形下二者彼此等价）。此外，按此对应，两侧对象的重要不变量（伽罗瓦一侧的 Frobenius 本征值，另一侧的 Hecke 本征值）应当相配。详见综述 [Frenkel 2007] 第一部分。

整体 Langlands 对应对于 $GL_2$ 已于 20 世纪 80 年代由 V. Drinfeld 证明（Drinfeld 1980; Drinfeld 1983; Drinfeld 1987; Drinfeld 1988）；最近 L. Lafforgue [2002] 对任意 $n$ 证明了 $GL_n$ 的情形。

整体对应与局部对应以下述意义相容。我们可以把每个局部域 $F_x$ 的 Weil 群 $W_{F_x}$ 嵌入整体 Weil 群 $W_F$。这样的嵌入不唯一，但差 $W_F$ 中的共轭是良定义的。因此 $W_F$ 的 $n$ 维表示的一个等价类 $\sigma$ 给出 $W_{F_x}$ 的 $n$ 维表示的一个良定义的等价类 $\sigma_x$，对每个 $x\in X$。由局部 Langlands 对应，可给 $\sigma_x$ 附上 $GL_n(F_x)$ 的不可约光滑表示的一个等价类。在此等价类中选一个表示 $\pi_x$。则对应于 $\sigma$ 的 $GL_n(\mathbb{A})$ 的自守表示同构于限制张量积 $\bigotimes'_{x\in X}\pi_x$。

% [Tutor Note: 含糊点拨（p.7）。作者说“先验地不清楚该张量积为何能实现在商上的函数空间里”，省略的正是整体对应的构造：限制张量积只给出一个光滑 $GL_n(\mathbb{A})$-模，证明它是自守的（即出现在合理的函数空间中）才是整体 Langlands 对应的内容。函数域情形由 Drinfeld（$n=2$）与 Lafforgue（任意 $n$）解决，细节见第 10.5 节。]

这是一个非常不平凡的陈述，因为先验地并不清楚为什么该张量积能实现在商 $GL_n(F)\backslash GL_n(\mathbb{A})$ 上的函数空间中。

% -------------------------- 印刷 p.8（PDF p.22） --------------------------

与局部情形一样，我们可能还想把群 $GL_n$ 换成 $F$ 上任意的约化代数群。一般的整体 Langlands 猜想粗略地预测：$G(\mathbb{A})$ 的不可约自守表示与同态 $W_F\to{}^L G$ 相关。但如局部情形，对一般约化群，猜想的精确表述要复杂得多（见 [Arthur 1989]）。

最后，整体 Langlands 猜想也可以对数域陈述，它们事实上正起源于此。此时取域 $F$ 为有理数域 $\Q$ 的有限扩张。例如考虑 $F=\Q$ 本身。已知 $\Q$ 的完备化（在同构意义下）是：对所有素数 $p$ 的 $p$-adic 数域 $\Q_p$（非阿基米德），以及实数域 $\R$（阿基米德）。于是素数扮演了有限域上代数曲线的点的角色（而阿基米德完备化在某种意义上对应于一个“无穷远点”）。阿代尔环 $\mathbb{A}_\Q$ 的定义与函数域情形相同，于是可以定义 $GL_n(\mathbb{A}_\Q)$ 或更一般约化群的自守表示的概念。猜测地，对伽罗瓦群 $\operatorname{Gal}(\overline{\Q}/\Q)$ 的 $n$ 维表示的每个等价类，都可附上 $GL_n(\mathbb{A}_\Q)$ 的不可约自守表示的一个等价类；但这一对应并不预期是双射，因为数域情形已知某些自守表示不对应任何伽罗瓦表示。

数域情形下的 Langlands 猜想导致非常重要且出人意料的结果。确实，伽罗瓦群的许多有趣表示可以在“自然”中找到。例如，群 $\operatorname{Gal}(\overline{\Q}/\Q)$ 会作用在定义于 $\Q$ 上的代数簇的几何不变量（例如 étale 上同调）上。于是，若取 $\Q$ 上的椭圆曲线 $E$，就会在其第一 étale 上同调上得到一个二维伽罗瓦表示。该表示包含曲线 $E$ 的许多重要信息，例如 $E$ 在不同素数 $p$ 下在 $\Z/p\Z$ 上的点数。Langlands 对应应当把这些伽罗瓦表示与 $GL_2(\mathbb{A}_F)$ 的自守表示联系起来，使得伽罗瓦一侧的数据（例如 $E(\Z/p\Z)$ 的点数）被翻译为自守一侧更易处理的东西，例如封装 $GL_2(\mathbb{A}_\Q)$ 的自守表示的模形式的 $q$-展开系数。\transnote{此句讨论 $\Q$ 上的椭圆曲线，却把自守表示写成 $GL_2(\mathbb{A}_F)$；本节前段讨论的是函数域 $\F_q(X)$，而此处域已取为 $\Q$，阿代尔环应是 $\mathbb{A}_\Q$——紧接的下一句正写作 $GL_2(\mathbb{A}_\Q)$。译文按原书保留 $\mathbb{A}_F$，公式未改。[依据：原书 p.8 同段上下文。][来源：本书 p.8。]}这导致若干惊人的推论，例如 Taniyama--Shimura 猜想。关于这一点，参见 [Frenkel 2007] 及其中参考文献。

% [Tutor Note: 原书笔误（p.8）。此段讨论 $E$ 定义在 $\Q$ 上，对应的应是 $GL_2(\mathbb{A}_\Q)$ 的自守表示；原文此处写作 $GL_2(\mathbb{A}_F)$，疑为 $GL_2(\mathbb{A}_\Q)$ 之误（紧接的下一句又写作 $GL_2(\mathbb{A}_\Q)$）。译文按原书保留，已登记报告。]

在函数域情形，Langlands 对应被证明更易研究。主要原因是函数域情形可以利用底曲线以及与该曲线相关的各种模空间的几何。曲线也可以在复数域上考虑。一些最近的结果表明，这样的曲线也存在整体 Langlands 对应的一个版本。这一对应的局部对应物，正是本书的主题。

% [Tutor Note: 背景（p.8）。此处指模性定理（原 Taniyama--Shimura 猜想，由 Wiles、Taylor 等人证明）：$\Q$ 上的椭圆曲线 $E$ 对应 $GL_2(\mathbb{A}_\Q)$ 的尖点自守表示（权 2 模形式），使 $a_p$ 与 $\#E(\Z/p\Z)$ 相配。它既是 Langlands 对应最著名的实例，也是 Fermat 大定理证明的关键一步；几何 Langlands 对应的许多猜想以它为原型。]

% 拆页衔接说明：本段末句“Some recent results ...”原文跨 PDF p.22–p.23。
% 该句的收尾（p.23 首行“show that a version ... for such curves”）已在此译出；
% U02 从 p.23 第二句“The local counterpart of this correspondence is the subject of this book.”起译，避免重复。
